\chapter{Queues}\label{ch:queues}

\section{The Queue ADT}

\begin{definition}[Queue]
A \emph{queue} is a sequence in which insertions take place at one end, the
\emph{rear} (or back), and deletions at the other, the \emph{front}. Its operations are
\cc{enqueue(x)} (also \cc{push}, \cc{add}, \cc{offer}), \cc{dequeue()} (also \cc{pop},
\cc{remove}, \cc{poll}), \cc{front()} (\cc{peek}), \cc{empty()} and \cc{size()}. The
item removed is always the one that has waited longest: \emph{first in, first out}
(FIFO).
\end{definition}

Queues model waiting lines: print jobs, requests to a server, processes waiting for the
processor (round-robin scheduling), messages between programs, and the frontier of a
breadth-first search. In each case fairness means serving in order of arrival.

\section{Array implementations}

Both ends of a queue change, so the array list of Chapter~\ref{ch:arrays}, which is fast
only at its rear, is not enough. There are three attempts.

\para{1. Front at index 0} Enqueue is \cc{items[size++] = x}, $\Theta(1)$; dequeue reads
\cc{items[0]} and shifts the remaining $n-1$ items left, $\Theta(n)$
(Theorem~\ref{thm:shift}).

\para{2. A moving front} Keep an index \cc{front} and dequeue by \cc{front++}. Both
operations become $\Theta(1)$, but the occupied part \cc{items[front..rear]} drifts to
the right. When it reaches the end of the array the queue reports ``full'' although the
cells before \cc{front} are free (\emph{false overflow}).

\para{3. The circular buffer} Let the indices wrap around: the cell after \cc{CAP-1} is
\cc{0}. With \cc{front} and \cc{size}, item $i$ of the queue is
\cc{items[(front+i) \% CAP]}; the rear item is $i=\cc{size}-1$ and the next free cell is
\cc{(front+size) \% CAP}.
\begin{itemize}
\item \cc{enqueue(x)}: \cc{items[(front+size)\%CAP] = x; size++;}
\item \cc{dequeue()}: \cc{x = items[front]; front = (front+1)\%CAP; size--;}
\end{itemize}
Every operation is $\Theta(1)$, and all \cc{CAP} cells are usable.

\listingfile[firstline=33,label=lst:queue]{lib/stackqueue.h}{The circular queue used in
Parts II and III (second part of \cc{stackqueue.h}).}

\begin{figure}[ht]\centering
\begin{tikzpicture}[x=7.5mm]
\foreach \v/\s [count=\i from 0] in {9/sh,4/sh,{}/,{}/,{}/,{}/,5/sh,7/sh} \node[cell,\s] (a\i) at (\i,0) {\v};
\foreach \i in {0,...,7} \node[idx,above=0pt of a\i] {\i};
\draw[ptr] (6,-0.9) node[below,lbl]{front=6} -- (a6.south);
\draw[ptr] (2,-0.9) node[below,lbl]{(front+size)\%8 = 2} -- (a2.south);
\node[right,font=\scriptsize,align=left] at (8.6,0) {CAP 8, size 4\\ queue (front to rear): 5 7 9 4};
\end{tikzpicture}
\caption{A circular queue that has wrapped around. The next \cc{enqueue} writes cell 2.}
\end{figure}

\begin{note}{Full or empty?}
Some textbooks store \cc{front} and \cc{rear} instead of \cc{size}. Then
\cc{front == rear} happens both when the queue is empty and when it is full. The usual
cure is to leave one cell always unused: the queue is empty when \cc{front == rear} and
full when \cc{(rear+1)\%CAP == front}, so the capacity is \cc{CAP-1}. Storing the size
(as here) avoids the problem and uses every cell. Read which convention a question uses
before you compute ``is it full?''.
\end{note}

\para{Growing} When the circular buffer is full, a growable version allocates an array
of twice the size and copies item $i$ to cell $i$ of the new array (not to the same
index!), so that the wrapped part is straightened; then \cc{front = 0}. As with the
vector, the cost is $\Theta(1)$ amortised per operation.

\section{Linked implementations}

In a singly linked list with head and tail, removal is $\Theta(1)$ only at the head and
insertion is $\Theta(1)$ at both ends (Chapter~\ref{ch:linked}). Hence the \emph{front}
of a linked queue is the \emph{head} and the \emph{rear} is the \emph{tail}: enqueue is
\cc{pushBack}, dequeue is \cc{popFront}. The opposite orientation would need
\cc{popBack}, which is $\Theta(n)$. Two details matter:
\begin{itemize}
\item dequeueing the last item must also set \cc{tail = NULL}, otherwise the next
  enqueue writes through a dangling pointer;
\item enqueueing into an empty queue must set \emph{both} \cc{head} and \cc{tail}.
\end{itemize}
The circular list of Section~\ref{sec:circular} gives a queue with a single pointer
\cc{last} (Listing~\ref{lst:circular}, Tracing~\ref{trace:24}).

\begin{table}[ht]\centering\small
\begin{tabular}{@{}lccc@{}}\toprule
 & front at index 0 & circular buffer & linked (head = front)\\\midrule
\cc{enqueue} & $\Theta(1)$ & $\Theta(1)$ (amortised if growable) & $\Theta(1)$\\
\cc{dequeue} & $\Theta(n)$ & $\Theta(1)$ & $\Theta(1)$\\
\cc{front}, \cc{size} & $\Theta(1)$ & $\Theta(1)$ & $\Theta(1)$\\\bottomrule
\end{tabular}
\caption{Queue implementations.}
\end{table}

\section{Queues built from stacks, and stacks from queues}\label{sec:conversions}

\para{A queue from two stacks} Push every new item on a stack \cc{in}. To dequeue, take
the top of a second stack \cc{out}; if \cc{out} is empty, first pop \emph{all} of
\cc{in} and push the items onto \cc{out}. The transfer reverses the order, so the oldest
item ends on top of \cc{out}.

\listingfile[label=lst:twostack]{sol/p43.c}{A queue from two stacks
(Problem~\ref{prob:43}, Tracing~\ref{trace:12}).}

\begin{theorem}[Two-stack queue]\label{thm:twostack}
Starting from an empty queue, any sequence of $m$ operations of
Listing~\ref{lst:twostack} performs at most $4m$ stack operations. Hence each queue
operation costs $\Theta(1)$ amortised, although a single dequeue may cost $\Theta(n)$.
\end{theorem}
\begin{proof}
Follow one item. It is pushed onto \cc{in} once, popped from \cc{in} once, pushed onto
\cc{out} once and popped from \cc{out} once: at most four stack operations, because
items never go back from \cc{out} to \cc{in}. The items are at most $m$ in number, and
besides moving items each queue operation does $O(1)$ work (one failed pop, one test),
so the total is $O(m)$.
\end{proof}

The order is correct because \cc{shift} moves items \emph{only when \cc{out} is empty}.
Moving whenever \cc{in} is non-empty would place newer items on top of older ones still
in \cc{out}, and the next dequeue would return a newer item first.

\para{A stack from one queue} After each enqueue of $x$, rotate the $n-1$ older items
(dequeue and re-enqueue each) so that $x$ is at the front. Then the front is always the
top: push is $\Theta(n)$, pop and top are $\Theta(1)$ (Problem~\ref{prob:44}). The
other choice --- push $\Theta(1)$, and pop rotates $n-1$ items to bring the newest to the
front --- costs the same in total.

\section{Applications}

\subsection{Rotation and the Josephus problem}

Dequeueing an item and enqueueing it again moves it from the front to the rear:
\emph{a queue is a circle}. In the Josephus problem $n$ people stand in a circle,
numbered $1..n$; counting starts at person 1, and every $k$-th person leaves the circle.
Which one remains? Simulate: rotate $k-1$ people, then remove the next.

\listingfile[label=lst:josephus]{sol/p45.c}{Josephus by simulation, $\Theta(nk)$, and by
the recurrence, $\Theta(n)$ (Problem~\ref{prob:45}).}

\begin{proposition}[Josephus recurrence]
Number the positions $0..i-1$, starting at the person where counting begins, and let
$J(i)$ be the position of the survivor among $i$ people. Then $J(1)=0$ and
$J(i)=\bigl(J(i-1)+k\bigr)\bmod i$ for $i\ge2$.
\end{proposition}
\begin{proof}
With $i$ people the first person removed is at position $(k-1)\bmod i$. Counting
restarts at the next person, position $k\bmod i$; renumber the remaining $i-1$ people
from there. The person with new number $p$ had old number $(p+k)\bmod i$. The survivor
has new number $J(i-1)$, hence old number $(J(i-1)+k)\bmod i$.
\end{proof}
For $n=7$, $k=3$ the removal order is $3,6,2,7,5,1$ and the survivor is $4$; the
recurrence gives $J=0,1,1,0,3,0,3$ for $i=1..7$, and $J(7)+1=4$.

\subsection{Sliding windows over a stream}

When items arrive in time order and only those of the last $w$ time units matter, the
items that expire are always the \emph{oldest}: they leave from the front, new ones join
at the rear. The queue holds exactly the current window (\emph{recent counter},
Problem~\ref{prob:47}). For a window of $k$ consecutive array positions, a queue of
\emph{indices} of interesting items suffices; an index leaves when it falls out of the
window (\emph{first negative in each window}, Problem~\ref{prob:49}). Each index enters
and leaves once, so the total cost is $\Theta(n)$, as for the monotonic stack.

When the question is the \emph{maximum} or \emph{minimum} of each window, items must
also leave from the \emph{rear}: this needs a deque (Section~\ref{sec:monodeque}).

\subsection{Combining a queue with a stack}

A stack reverses, a queue preserves order. Pouring part of a queue into a stack and back
reverses that part. Reversing the first $k$ items of a queue
(Problem~\ref{prob:46}): dequeue $k$ items onto a stack, pop them back into the queue
(now they are reversed, but at the rear), then rotate the remaining $n-k$ items once
around. Interleaving the two halves of a queue (Problem~\ref{prob:48}) is solved the
same way with a few more passes. In such problems, remember that \cc{qSize(\&q)}
\emph{changes} while the loop runs: read it into a variable first.

\begin{table}[ht]\centering\small
\begin{tabular}{@{}p{0.42\linewidth}p{0.51\linewidth}@{}}\toprule
If the problem says \dots & think of \dots\\\midrule
in order of arrival, first come first served, round robin & a plain queue\\
every $k$-th, circle, rotate & dequeue + enqueue (\ref{prob:45})\\
last $w$ milliseconds, recent, expire & queue of time stamps (\ref{prob:47})\\
something about each window of size $k$ & queue of indices (\ref{prob:49}); deque if min/max (\ref{prob:53})\\
reverse part of a queue & a stack as helper (\ref{prob:46})\\
implement X using only Y & Section~\ref{sec:conversions} (\ref{prob:43}, \ref{prob:44})\\\bottomrule
\end{tabular}
\caption{Recognising a queue problem.}
\end{table}
